% Simulazione d'esame: prova del 01/02/1970, esercizio 3 -- shadowgraph di un ugello a spina troncata.
% Stessi numeri e figure dello svolgimento di riferimento:
% svolgimenti/1970-02-01_E03_ugello_spina_troncata (calcoli.py e figure/).
% Compilare con: bash ../compila.sh 1970-02-01_E03_ugello_spina_troncata  (A4 + orologio)

% Preambolo comune delle simulazioni d'esame (vedi GUIDA.md).
% Lo stesso sorgente si compila in due formati:
%   - A4 (svolgimento.pdf), da leggere su carta, PC o telefono;
%   - orologio (orologio.pdf): colonna stretta, letta dall'app dell'orologio
%     (garmin-latex), che taglia le pagine in strisce larghe quanto lo schermo.
% Il formato orologio si sceglie definendo \formatoorologio prima di \input
% (lo fa compila.sh). In quel formato il margine destro è volutamente largo:
% ciò che sporge oltre la colonna (equazioni lunghe, tabelle, figure) viene
% riconosciuto dall'app e mostrato rimpicciolito, con lo zoom.

\documentclass[11pt]{article}

\usepackage[T1]{fontenc}
\usepackage[utf8]{inputenc}
\usepackage[italian]{babel}
\usepackage{amsmath,amssymb}
\usepackage{array}
\usepackage{siunitx}
\usepackage{booktabs}
\usepackage{graphicx}
\usepackage{float}
\usepackage{caption}
\usepackage{enumitem}
\usepackage{ragged2e}
\usepackage{titlesec}
\usepackage{hyperref}

\hypersetup{hidelinks}
\setlength{\parindent}{0pt}
\setlength{\parskip}{4pt}
\allowdisplaybreaks
\sisetup{output-decimal-marker={,}, group-separator={\,}, per-mode=symbol,
         exponent-product=\cdot, range-phrase={\text{--}}}

% Notazione del formulario (come template/svolgimento.tex).
\newcommand{\Cp}{C^{+}}
\newcommand{\Cm}{C^{-}}
\newcommand{\Mr}{M_{r}}
\newcommand{\Ms}{M_{s}}

% Separatore fra relazioni nella stessa formula: \begin{gather*} A \ee B \end{gather*}.
% In A4 le relazioni stanno sulla stessa riga, nel formato orologio una sotto l'altra.
\ifdefined\formatoorologio \newcommand{\ee}{,\\} \else \newcommand{\ee}{,\qquad} \fi

% Titolo dell'esercizio: \esercizio{Prova del GG/MM/AAAA -- Esercizio n}{argomento}.
% Va anche nei metadati del PDF: l'app dell'orologio lo usa come titolo nell'elenco.
\newcommand{\esercizio}[2]{%
  \title{#1\\[2pt]\large #2}\author{}\date{}%
  \hypersetup{pdftitle={#1: #2}}}

\ifdefined\formatoorologio
  % Colonna di 4,6 cm: a 2,4 px/pt è larga 262 px sullo schermo da 390.
  \usepackage[paperwidth=12cm,paperheight=150cm,left=0.2cm,top=0.2cm,
              textwidth=4.6cm,textheight=149cm,nohead,nofoot]{geometry}
  \usepackage[scaled=1.0]{helvet}
  \usepackage{sfmath}
  \renewcommand{\familydefault}{\sfdefault}
  \pagestyle{empty}
  \makeatletter\let\ps@plain\ps@empty\makeatother   % anche la pagina del titolo
  % colonna dichiarata per l'app dell'orologio (deve coincidere con left e textwidth)
  \hypersetup{pdfkeywords={colonna=0.2cm+4.6cm}}
  \RaggedRight
  \setlength{\parskip}{5pt}
  \hfuzz=20cm          % ciò che sporge va nel margine destro, senza avvisi
  \hbadness=10000
  \titleformat*{\section}{\normalsize\bfseries\RaggedRight\hyphenpenalty=10000}
  \titleformat*{\subsection}{\normalsize\bfseries\RaggedRight}
  \setlist{leftmargin=*,itemsep=2pt}
  \captionsetup{font=small,justification=raggedright,singlelinecheck=false}
  % figure: larghe il doppio della colonna, così sporgono e l'app le mostra con lo zoom
  \newcommand{\figura}[3][0.7]{%
    \begin{figure}[H]\includegraphics[width=9cm]{#2}\caption{#3}\end{figure}}
\else
  \usepackage[a4paper,margin=2.5cm]{geometry}
  \newcommand{\figura}[3][0.7]{%
    \begin{figure}[H]\centering\includegraphics[width=#1\textwidth]{#2}\caption{#3}\end{figure}}
\fi

\graphicspath{{../../svolgimenti/1970-02-01_E03_ugello_spina_troncata/figure/}}

\esercizio{Prova del 01/02/1970 -- Esercizio 3}{Shadowgraph di un ugello a spina troncata}

\begin{document}
\maketitle

\emph{Traccia.} Identificare le onde e le regioni numerate nello shadowgraph della corrente supersonica stazionaria
generata da un ugello a spina (plug) troncato bidimensionale (4/30). Accanto c'è lo schema della geometria (metà
superiore). È la stessa traccia della prova del 25/01/2023.

\section*{Che cosa mostra la figura}
Dallo schema riconosco la metà superiore dell'ugello: una carenatura con il bordo interno curvo, la gola fra il
labbro della carenatura e la sommità della spina, la rampa della spina che scende verso l'asse e il troncamento, una
base piana perpendicolare all'asse. L'ugello è piano e simmetrico: lo shadowgraph mostra entrambe le metà, con i due
getti che escono dalle fessure (labbri in alto e in basso a sinistra), il corpo nero della spina al centro e la base
sulla linea punteggiata verticale. La corrente va da sinistra verso destra.

Lo shadowgraph è sensibile alle variazioni del gradiente di densità: un urto appare come una coppia di linee
scura--chiara molto netta, uno strato di scorrimento come una linea più sfumata, la testa o la coda di un ventaglio
d'espansione, dove il gradiente di densità cambia bruscamente, come una linea singola. Le strutture sono nitide e
simmetriche: il campo è stazionario, come dice il testo. Non ho scale né pressioni, quindi riferisco le lunghezze
alla semialtezza della base $h_b$ e le pressioni alla pressione totale in camera $p_0$, e assumo aria,
$\gamma=1{,}4$. Il reticolo fine (foto dello schermo), i puntini del contorno della spina, la cornice della finestra
e le linee bianche dei rimandi ai numeri sono artefatti.

\figura[0.95]{annotata.png}{Shadowgraph (raddrizzato) con le strutture che individuo. Arancio: rampe e base; ciano
(1): contorni esterni del getto; giallo (2): linea allo spigolo della base; verde (3): contorni interni del getto;
rosso (4): urti obliqui di ricompressione; magenta: punto di ristagno al collo (5) e regione di base (6). Le lettere
$r$, $c$, $j$, 2 sono gli stati della metà superiore.\label{fig:ann}}

\section*{Modello e misure}
Gli urti obliqui netti (4) e i ventagli centrati negli spigoli esistono solo in corrente supersonica: nei getti uso
quindi il modello 2D stazionario supersonico (ventagli di Prandtl--Meyer, urti obliqui, caratteristiche $\Cp$ e
$\Cm$), mentre la regione di base la tratto come una regione a pressione uniforme $p_b$. Assumo la corrente
isentropica dalla camera fino all'urto 4, con $p_0$ uniforme sulla linea di corrente che costeggia il contorno interno
(trascuro il lip shock, che è debole). Sui contorni del getto impongo le condizioni giuste: sul contorno esterno
$p=p_a$, su quello interno $p=p_b$ vicino allo spigolo e la direzione della corrente uguale a quella del contorno,
che misuro. Calcolo la metà superiore e, per gli angoli, uso le medie delle due metà, che differiscono di pochi gradi.

Con righello e goniometro, prendendo l'asse come la mezzeria fra i due spigoli della base (angoli dall'asse, verso
l'asse se non dico altro), misuro:
\begin{itemize}
\item rampa della spina nel tratto finale \ang{26} (più ripida a monte, circa \ang{45});
\item contorno interno (3): \ang{37} subito dopo lo spigolo, \ang{32} a metà, \ang{25} prima del collo;
\item linea allo spigolo (2): circa \ang{5};
\item urto obliquo (4): \ang{21} sull'asse, verso l'esterno;
\item contorno esterno (1): circa \ang{3};
\item collo (5) a $L\approx1{,}9\,h_b$ dalla base; gli urti 4 arrivano sui contorni esterni a circa $6\,h_b$.
\end{itemize}

\section*{Il getto e il suo contorno esterno (1)}
Dal labbro della carenatura parte il \textbf{contorno esterno del getto} (1): è una superficie libera, cioè uno strato
di scorrimento che separa il getto dall'aria ambiente in quiete, con pressione $p_a$ e velocità tangente; per
viscosità diventa uno strato di mescolamento che si allarga verso valle. È quasi rettilineo e quasi parallelo
all'asse (\ang{3}). Fra il labbro e la rampa vedo le striature del ventaglio d'espansione centrato nel labbro (onde
$\Cm$), che porta la corrente esterna alla pressione ambiente: è l'espansione esterna della spina, quella che le
permette di adattarsi alla quota.

\figura[0.95]{schema.pdf}{Schema della metà superiore in unità di $h_b$, con le strutture 1--6, i ventagli (blu), le
onde di compressione del contorno interno (rosso sottile) e l'espansione in cui l'urto 4 si riflette sul contorno
esterno.\label{fig:schema}}

Per il resto seguo una linea di corrente vicina al contorno interno della metà superiore (fig.~\ref{fig:schema}),
con gli stati $r$ (fine rampa), $c$ (dopo lo spigolo, a $p_b$), $j$ (prima del collo) e 2 (dietro l'urto 4). Le
zone a due famiglie (onde del labbro riflesse dalla rampa, compressioni che attraversano il ventaglio dello spigolo)
non le risolvo: le tratto come raccordi fra gli stati uniformi.

\section*{Al collo: l'urto obliquo (4) e il punto di ristagno (5)}
Comincio dall'urto, perché è la struttura più netta (coppia scura a monte, chiara a valle) e dà la misura più
affidabile. I due \textbf{contorni interni} (3) si chiudono sull'asse nel \textbf{collo}: lì la linea di corrente che
separa il fluido del getto da quello che ricircola arriva sull'asse con velocità nulla, perché la ricompressione non
lascia proseguire il fluido lento dello strato di scorrimento, che viene rimandato verso la base. Il collo è quindi il
\textbf{punto di ristagno posteriore della scia chiusa} (5). Da qui partono i due \textbf{urti obliqui di
ricompressione} (4) e, fra di essi, la scia lungo l'asse. L'urto 4 è della famiglia $\Cp$ e raddrizza parallelamente
all'asse la corrente che convergeva verso di esso: le due metà si raccordano come in una riflessione regolare
sull'asse.

La corrente a monte ha la direzione del contorno interno prima del collo, a valle è assiale: deviazione
$\theta=\ang{25}$. L'urto forma \ang{21} con l'asse dalla parte opposta, quindi con la corrente a monte
$\sigma=\ang{21}+\ang{25}=\ang{46}$. La relazione $\theta$--$\sigma$--$M$
\[
\tan\theta=2\cot\sigma\,\frac{M^2\sin^2\sigma-1}{M^2(\gamma+\cos2\sigma)+2}
\]
è lineare in $M^2$, quindi la inverto esplicitamente:
\[
M_j^2=\frac{2(\tan\theta+\cot\sigma)}{2\cot\sigma\sin^2\sigma-\tan\theta\,(\gamma+\cos2\sigma)}
\]
e trovo $M_j\approx2{,}8$. Controllo subito il ramo: $\theta=\ang{25}<\theta_{\max}(2{,}8)=\ang{32.7}$ e
$\sigma=\ang{46}$ è minore dell'angolo del punto di $\theta_{\max}$ (\ang{65}): ramo debole, come deve essere in campo
aperto.

Attraverso l'urto uso le relazioni dell'urto normale sulla componente normale:
\begin{gather*}
M_{n,j}=M_j\sin\sigma\approx2{,}02 \ee
\frac{p_2}{p_j}=1+\frac{2\gamma}{\gamma+1}\left(M_{n,j}^2-1\right)\approx4{,}6,
\end{gather*}
\begin{gather*}
M_{n2}^2=\frac{1+\delta M_{n,j}^2}{\gamma M_{n,j}^2-\delta} \ee
M_2=\frac{M_{n2}}{\sin(\sigma-\theta)}\approx1{,}6 .
\end{gather*}
Dietro l'urto la corrente resta supersonica ($M_2\approx1{,}6$), quindi il modello 2D stazionario vale anche a valle
del collo, fuori dalla scia; e $p_{02}/p_0\approx0{,}71<1$: l'entropia cresce e l'urto è di compressione
($p_2>p_j$), come deve.

\section*{Lungo il contorno interno (3): la regione di base (6) e la pressione di base}
Il \textbf{contorno interno del getto} (3) parte dallo spigolo della base e separa il getto dalla regione di base;
è curvo: scende verso l'asse sempre meno inclinato (da \ang{37} a \ang{25}), quindi ruota la corrente in senso
antiorario e la comprime con onde $\Cp$ che coalescono nell'urto 4. La regione che racchiude con la base fino al
collo è la \textbf{scia chiusa} (6), cioè la regione di base: il fluido vi ricircola lentamente (subsonico) a
pressione quasi uniforme $p_b$, la pressione di base. È chiusa perché è circondata da corrente supersonica e non
comunica con l'ambiente; fa da \textbf{poppa fluida}, cioè sostituisce la parte di spina tolta con il troncamento.

Il ventaglio dello spigolo e le compressioni del contorno curvo sono onde $\Cp$ (nascono dal contorno, che sta sotto
la corrente della metà superiore): attraverso di esse si conserva l'invariante delle $\Cm$,
\begin{gather*}
\theta+\nu=\text{cost} \ee
\nu_c=\nu_j+\theta_j-\theta_c \ee
\nu_r=\nu_j+\theta_j-\theta_r .
\end{gather*}
Con $\nu_j=\nu(M_j)\approx\ang{45.9}$, $\theta_j=\ang{-25}$, $\theta_c=\ang{-37}$, $\theta_r=\ang{-26}$:
\begin{gather*}
\nu_c\approx\ang{57.9},\ M_c\approx3{,}5 \ee
\nu_r\approx\ang{46.9},\ M_r\approx2{,}9 .
\end{gather*}
Dall'isentropica, $p/p_0=(1+\delta M^2)^{-\gamma/(\gamma-1)}$:
\begin{gather*}
\frac{p_r}{p_0}\approx0{,}034 \ee
\frac{p_b}{p_0}=\frac{p_c}{p_0}\approx0{,}014 \ee
\frac{p_j}{p_0}\approx0{,}036 \ee
\frac{p_2}{p_0}=\frac{p_2}{p_j}\frac{p_j}{p_0}\approx0{,}17 .
\end{gather*}
La pressione di base è circa il 40\% di quella di fine rampa ($p_b/p_r\approx0{,}41$) e la ricompressione al collo la
moltiplica per circa 12. Controllo la natura delle onde: allo spigolo $\theta$ diminuisce (da \ang{-26} a \ang{-37})
attraverso onde $\Cp$, quindi è un'espansione e $p_b<p_r$; lungo il contorno interno $\theta$ aumenta, quindi è una
compressione, coerente con l'urto 4 che ne è la coalescenza. Controllo anche la geometria della scia: l'inclinazione
media del contorno fra spigolo e collo, $\arctan(h_b/L)\approx\ang{28}$, sta fra quella iniziale (\ang{37}) e quella
finale (\ang{25}).

\section*{Allo spigolo della base (2)}
Nello spigolo la corrente della rampa gira verso l'asse per adattarsi a $p_b$: nasce un \textbf{ventaglio
d'espansione centrato} (onde $\Cp$) e, a ridosso dello spigolo, nello strato limite che si stacca, un urto debole di
ricompressione, il \textbf{lip shock}. Per capire che cosa sia la linea 2, corta e un po' curva, calcolo la testa e la
coda del ventaglio, inclinate di $\theta+\mu$ con $\mu=\arcsin(1/M)$:
\begin{gather*}
\theta_r+\mu_r\approx\ang{-26}+\ang{20.5}=\ang{-5.5} \ee
\theta_c+\mu_c\approx\ang{-37}+\ang{16.8}=\ang{-20} .
\end{gather*}
La linea 2 è inclinata di circa \ang{-5}: è la testa del ventaglio, e il lip shock non si distingue da essa. La coda
non si vede perché le compressioni del contorno curvo la attraversano subito.

\section*{Più a valle}
L'urto 4 arriva sul contorno esterno a circa $6\,h_b$ dalla base; lì la pressione è imposta ($p_a$), quindi si
riflette come espansione $\Cm$ e il contorno piega verso l'esterno. Nella figura si intravedono, più a valle, le onde
riflesse che si incrociano vicino all'asse: è l'inizio della struttura a celle del getto.

\section*{Risposta}
\begin{enumerate}
\item \textbf{Contorno esterno del getto}: superficie libera (strato di scorrimento) a $p=p_a$, dal labbro della
  carenatura.
\item \textbf{Struttura allo spigolo della base}: ventaglio d'espansione $\Cp$ centrato nello spigolo (se ne vede la
  testa), con il lip shock a ridosso dello spigolo.
\item \textbf{Contorno interno del getto}: strato di scorrimento curvo fra il getto e la regione di base, a $p_b$
  vicino allo spigolo.
\item \textbf{Urto obliquo di ricompressione}: urto $\Cp$ dal collo, raddrizza la corrente lungo l'asse.
\item \textbf{Punto di ristagno al collo}: chiusura della scia, dove i contorni interni si incontrano sull'asse.
\item \textbf{Scia chiusa} (regione di base): ricircolazione subsonica a pressione uniforme $p_b$, poppa fluida.
\end{enumerate}

\begin{table}[H]
  \centering
  \begin{tabular}{cccc}
    \toprule
    Stato & $\theta$ & $M$ & $p/p_0$ \\
    \midrule
    $r$ & \ang{-26} & $\approx2{,}9$ & $\approx0{,}034$ \\
    $c$ & \ang{-37} & $\approx3{,}5$ & $\approx0{,}014$ \\
    $j$ & \ang{-25} & $\approx2{,}8$ & $\approx0{,}036$ \\
    2 & \ang{0} & $\approx1{,}6$ & $\approx0{,}17$ \\
    \bottomrule
  \end{tabular}
  \caption{Stati lungo il contorno interno (metà superiore); $p_c=p_b$.}
\end{table}

\figura[0.98]{p_theta.pdf}{Piano $p$--$\theta$ (pressioni riferite a $p_0$), in scala logaritmica e lineare: curva
isentropica $\Cp$ per $r$, $c$, $j$ (blu), polare $\Cp$ di $j$ (rosso; sottile il ramo forte), le stesse curve
speculari per la metà inferiore (punteggiate) e la pressione di base $p_b$.\label{fig:pth}}

Nel piano $p$--$\theta$ (fig.~\ref{fig:pth}) gli stati $r$, $c$, $j$ stanno sulla stessa curva isentropica, perché
il ventaglio dello spigolo e le compressioni del contorno interno sono onde della stessa famiglia $\Cp$: la corrente
scende da $r$ a $c$ (fino a $p_b$, la retta tratteggiata) e risale fino a $j$. Il ramo debole della polare $\Cp$ di
$j$ porta a $\theta=0$ nello stato 2, dove passa anche la polare speculare della metà inferiore: stessa pressione e
stessa direzione sull'asse, come in una riflessione regolare. Le pressioni coprono un ordine di grandezza: la scala
logaritmica mostra bene la zona di bassa pressione intorno a $r$, $c$, $j$, che quella lineare schiaccia sull'asse
delle ascisse; quella lineare dà la proporzione reale del salto all'urto, molto maggiore delle variazioni lungo il
contorno interno.

\section*{Discussione}
\begin{itemize}
\item Il modello ritrova ciò che si vede: la testa del ventaglio con l'inclinazione della linea 2, il contorno
  interno curvo con le compressioni che coalescono nell'urto 4, l'urto sul ramo debole con corrente supersonica
  parallela all'asse a valle, le onde riflesse dai contorni esterni.
\item \textbf{Scia chiusa e pressione di base.} La regione di base è racchiusa da corrente supersonica, quindi
  l'ambiente non vi fa arrivare informazione: $p_b/p_0$ non dipende da $p_a$ e resta fissata dal getto (qui circa
  $0{,}014$). Il contorno esterno invece è a $p_a$ e si adatta alla quota: è il vantaggio della spina rispetto a un
  ugello a campana. A rapporti di pressione più bassi la scia è aperta: la regione di base comunica con l'ambiente
  attraverso una scia subsonica e $p_b$ segue $p_a$.
\item \textbf{Poppa fluida.} Il contorno interno, curvo da \ang{37} a \ang{25}, prende il posto della parte di rampa
  tolta, che avrebbe portato la corrente in direzione assiale con un'espansione senza urti. Con il troncamento la
  ricompressione finale avviene invece con l'urto 4, con una perdita di pressione totale (circa il 30\% sulla linea di
  corrente vicina al contorno interno) e una base a pressione minore di quella di fine rampa, che sulla base dà una
  spinta ridotta.
\item \textbf{Limiti del modello.} Gli strati di scorrimento hanno spessore finito e crescono verso valle (il collo è
  una zona di dimensione finita, non un punto); il lip shock nasce nello strato viscoso e il modello non viscoso non lo
  descrive; la ricompressione al collo è distribuita fra onde e urto. Per $M$ elevati $\nu(M)$ varia poco, quindi
  $M_c\approx3{,}5$ e $p_b$ sono solo indicativi; $M_j$, ricavato direttamente dall'urto, è la stima più robusta.
\end{itemize}

\end{document}
